\scanpage{122}

proposes to examine the ways of \deleted{extending}\inserted{lifting}
the action of $\Gamma$ on $\mathcal X$ \deleted{(via $M$)} to an
action \deleted{of $\mathcal M$}\inserted{on $\mathcal X'_0$}, compatible
with the action \inserted{of $\Gamma$ on the group of operators $M$}
\deleted{on $\mathcal X$ via $\Gamma$,
\inserted{and extend the action of
$M\subset\mathfrak S$,}
\unsure{or what} \unsure{amounts} to the same thing, by an action of
$\mathfrak S$, compatible with its action on
$\mathcal X$ (via $\Gamma$)}.
More generally, if $\mathcal X'$ is an $M_{\mathcal X}$-torsor on
$\mathcal X$ isomorphic \unsure{to} $\mathcal X'_0$, hence obtained
\deleted{\unsure{by a}}\inserted{from $\mathcal X'_0$} by twisting with
an $M$-torsor $E$ ($\unsure{\simeq|=}$ fiber of $\mathcal X'_0$ at $x_0$),
let us study the analogous \deleted{extensions}\inserted{liftings} of the action
of $\deleted{\unsure{\mathfrak S}}\inserted{\Gamma}$ on $\mathcal X$ to an
action of $\deleted{\unsure{\mathfrak S}}\inserted{\Gamma\text{ on }\mathcal X'}$,
compatible with the action of
$\deleted{\unsure{\mathfrak S}}\inserted{\Gamma}$ on $\mathcal X$.

\deleted{Note that, if $\widetilde{\mathcal X}$ denotes the universal covering of
$\mathcal X$ at $x_0$, $\mathcal M$ operates on $\widetilde{\mathcal X}$ in a way
compatible with its action on $\mathcal X$, \unsure{while|and} extending the
\inserted{canonical} action of $\Pi$ on $\widetilde{\mathcal X}$ --- as
$[\Pi,\Pi]\subset\Pi$ is normal in $\mathcal M$, this action passes to the
quotient as a canonical action of
$\mathfrak S=\mathcal M/[\Pi,\Pi]$ on
$\mathcal X'_0=\widetilde{\mathcal X}/[\Pi,\Pi]$, which satisfies the}

\margin{121}

\begin{center}
\small
121\\
\texttt{https://grothendieck.umontpellier.fr}\\
Université de Montpellier, UMR CNRS 5149, Case courrier 051, Place Eugène Bataillon\\
34095 Montpellier cedex 5 - France
\end{center}

\scanpage{123}
\mspage{709}

\deleted{\inserted{two} conditions envisaged. Let $\Theta_0$ be this
operation of $\mathfrak S$, and $\Theta$ an arbitrary operation,
satisfying the said conditions.}

We thus want to make $\mathcal X'$, which is a covering of $\mathcal X$, given
by its fiber $E$ at $x_0$ and the action of $\Pi=\Pi_1(\mathcal X,x_0)$ on
$E$ (by translations, thanks to an $M$ torsor structure on $E$) into a
covering of $(\mathcal X',\deleted{\unsure{\mathfrak S}}\inserted{\Gamma})$
--- i.e. \uline{extend} the given action of $\Pi$ on \unsure{the|the} fiber
$E$, to an action of $\mathcal M$ on $E$ --- which moreover necessarily factors
through
$\deleted{\unsure{\mathfrak S}}\inserted{\widetilde{\mathcal M}}$.
I say that this action necessaril\textsuperscript{y} goes \unsure{by} way of the
affine group of the torsor $E$ --- to be precise, $\Gamma$ operates on $M$,
\unsure{let} $\gamma_M$ be the operation defined by $\gamma\in\Gamma$, and let
$g\in\deleted{\unsure{\mathfrak S}}\inserted{\widetilde{\mathcal M}}$ be over
$\gamma\in\Gamma$, we will thus have, for $x\in E$, $\lambda\in M$
\[
\tag{280}
\diagram{\begin{tikzcd}[column sep=0em,row sep=0.6em]
g & (\lambda & {}\cdot{} & x & {})=\gamma_M(\lambda)\mathbin{.}gx\\
\deleted{\unsure{\mathfrak S}}\inserted{\widetilde{\mathcal M}}
\arrow[u,phantom,"\in"{description,rotate=-90}]
& M\arrow[u,phantom,"\in"{description,rotate=-90}]
& & E\arrow[u,phantom,"\in"{description,rotate=-90}] &
\end{tikzcd}}
\]
which \unsure{is} trivial by calculation.
\deleted{\unsure{for|is}}
\[
g(\lambda.x)=(g\lambda)x=(g\lambda g^{-1}g)x
=\underbrace{(g\lambda g^{-1})}_{\gamma_M(\lambda)}.(gx)
\]

\margin{122}

\begin{center}
\small
122\\
\texttt{https://grothendieck.umontpellier.fr}\\
Université de Montpellier, UMR CNRS 5149, Case courrier 051, Place Eugène Bataillon\\
34095 Montpellier cedex 5 - France
\end{center}

\scanpage{124}

Thus if we introduce the group $\operatorname{Aff}(E)$, we will have a
homomorphism of exact sequences
\[
\tag{281}
\diagram{\begin{tikzcd}[column sep=large,row sep=large]
1 \arrow[r] & M \arrow[r,dash]
\arrow[d,Rightarrow,"{\deleted{\unsure{\Theta}}\mathrm{id}_M}"]
& \deleted{\unsure{\mathfrak S}}
\inserted{\widetilde{\mathcal M}} \arrow[r]
\arrow[d,"{\Theta_{\deleted{\unsure{\mathfrak S}}\inserted{\widetilde{\mathcal M}}}}"]
& \Gamma \arrow[r,dash] \arrow[d,"\Theta_\Gamma"] & 1\\
1 \arrow[r,dash] & M \arrow[r]
& \operatorname{Aff}(E) \arrow[r]
& \operatorname{Aut}(M) \arrow[r] & 1
\end{tikzcd}}
\]
where the outer vertical arrows $\mathrm{id}_M$ \inserted{and
$\Theta_\Gamma$} are known, and where the middle arrow is to be determined.
We can say that the category of the $\mathcal X'$ \unsure{sought|sought}
(with the operation of $\Gamma$ on them satisfying the desired conditions)
\inserted{is} equivalent to that of pairs
$(E,\Theta_{\deleted{\unsure{\mathfrak S}}\inserted{\widetilde{\mathcal M}}})$,
where $E$ is an $M$-torsor, and
$\Theta_{\deleted{\unsure{\mathfrak S}}\inserted{\widetilde{\mathcal M}}}$
a homomorphism of groups
$\deleted{\unsure{\mathfrak S}}\inserted{\widetilde{\mathcal M}}
\longrightarrow\operatorname{Aff}(E)$ \deleted{\unsure{making}}
compatible with $\mathrm{id}_M$ and $\Theta_\Gamma$ i.e. which
\unsure{makes|makes} (281) commutative.
We can also say: once $E$ is chosen,
\inserted{(and this choice is \unsure{harmless} up to nonunique
isomorphism\unsure{!!!|\ldots})}
that this defines an extension $\operatorname{Aff}(E)$ of
$\operatorname{Aut}(M)$ by $M$, whence thanks to
$\Gamma\longrightarrow\operatorname{Aut}(M)$ an
extension

\margin{123}

\begin{center}
\small
123\\
\texttt{https://grothendieck.umontpellier.fr}\\
Université de Montpellier, UMR CNRS 5149, Case courrier 051, Place Eugène Bataillon\\
34095 Montpellier cedex 5 - France
\end{center}

\scanpage{125}
\mspage{\textcircled{710}}

\margin{the splittings of $\operatorname{Aff}(E)$ determined by the points of $E$ -- \unsure{shift} of the origin \unsure{inner} \unsure{splittings}\ldots}

by $M$, \inserted{which is \uline{moreover split}.} (That said, the liftings of the action of $\Gamma$ on it correspond to the isomorphisms of this extension of $\Gamma$ by $M$, with the extension $\mathfrak S$ itself. Thus it exists iff the extension of $\mathfrak S$ of $\Gamma$ by $M$ is split, \deleted{\unsure{which}} \inserted{i.e.\ \unsure{the expression}}
\[
c^2(\mathfrak S)\in H^2(\Gamma,M)\quad\text{is zero.}
\tag{282}
\]
If such a lifting exists i.e. such an isomorphism, these form a torsor \unsure{under|over} the group of automorphisms of the extension (trivial) $\mathfrak S$ of $\Gamma$ by $M$ which is identified with the group
\[
(\qquad Z^1(\Gamma,M)
\]
of crossed homomorphisms from $\Gamma$ to $M$. The $M$-conjugacy classes of such \deleted{liftings \unsure{(or}} isomorphisms therefore correspond to the elements of
\[
H^1(\Gamma,M).
\]

Another way of saying it, by taking $E={}$trivial \inserted{$M$-}torsor${}=M$, so $\operatorname{Aff}(E)$ split, is that the chosen liftings correspond to the \uline{splittings} of the extension $\mathfrak S$ of $\Gamma$ by $M$, and the \inserted{$M$-}conjugacy classes of the liftings of the action of $\Gamma$ on $\mathcal X'_0$

\margin{124}

\begin{center}
\small
124\\
\texttt{https://grothendieck.umontpellier.fr}\\
Université de Montpellier, UMR CNRS 5149, Case courrier 051, Place Eugène Bataillon\\
34095 Montpellier cedex 5 - France
\end{center}

\scanpage{126}

correspond to the $M$-conjugacy classes of the splittings of the said extension $\mathfrak S$ of $\Gamma$ by $M$.

This being well set out, let us return to the situation of
\[
\left(\underbrace{U_{0,3\overline{\mathbb Q}}}_{\mathcal X},\quad
\Gamma=\underbrace{\Gamma_{\mathbb Q}\times\mathfrak S_3}_{N_{0,3}}\right),
\]
the $\Gamma$-coverings of $\mathcal X$ correspond to the $\mathfrak S_3$-coverings of $U_{0,3\mathbb Q}$, i.e. to the coverings of $(U_{0,3\mathbb Q},\mathfrak S_3)$.

\margin{\uline{N.B.} The group we have just denoted $M$ becomes here $\unsure{\widehat M}$, $\mathcal M$ becomes $\mathcal M_{0,3}$, $\pi$ becomes $\widehat\Pi_{0,3}$.}

Using the general constructions of \S~49 $\overline{\uline{\mathrm{VI}}}$, \deleted{we deduced} from the quotient
\[
\deleted{\unsure{G_{\rho,\sigma}}=\unsure{\mathfrak S}^{+\wedge}/\unsure{\mu_G}\simeq
\operatorname{SL}(2,\mathbb Z)^\wedge/
[\unsure{\widehat\Pi_{0,3}},\unsure{\widehat\Pi_{0,3}}]}
\]
\[
\mathfrak S^{+\wedge}/\mu_{\mathfrak S}\simeq
\unsure{\mathcal M}_{0,3}^{+}/[\widehat\Pi_{0,3},\widehat\Pi_{0,3}]
\subset\widetilde{\mathcal M}
\]
\deleted{is}

\margin{we could denote it $\unsure{\widetilde{\mathcal M}'}$!}

\deleted{ext} (with the notations of the preceding pages), extension of $\mathfrak S_3\subset\Gamma(=\Gamma_{\mathbb Q}\times\mathfrak S_3)$ by $M^\wedge$, we \inserted{deduced} a profinite group, \deleted{\unsure{and}} which was denoted $G_{\rho,\sigma}/\mu_G$ \inserted{or $G_{\rho,\widetilde\sigma}$} in the preceding sections and appeared there as a $\frac12$-direct product
\[
G_{\rho,\widetilde\sigma}
=\deleted{\unsure{\mathbb O|\Phi}}\widehat{\mathbb Z}^{*}\mathbin{.}\widetilde{\mathcal M},
\qquad\text{extension of }\widehat{\mathbb Z}^{*}\times\mathfrak S_3
\quad\text{by }M^\wedge
\]
and a canonical homomorphism
\[
\widetilde{\mathcal M}
=\mathcal M_{0,3}/(\widehat\Pi_{0,3},\widehat\Pi_{0,3})
\longrightarrow G_{\rho,\widetilde\sigma}
\]
fitting into a homomorphism of exact sequences

\margin{125}

\begin{center}
\small
125\\
\texttt{https://grothendieck.umontpellier.fr}\\
Université de Montpellier, UMR CNRS 5149, Case courrier 051, Place Eugène Bataillon\\
34095 Montpellier cedex 5 - France
\end{center}

\scanpage{127}
\mspage{711}

\diagram{
\[
\begin{tikzcd}[column sep=large,row sep=huge]
1 \arrow[r,no head] &
\widehat M \arrow[r] \arrow[d,Rightarrow] &
\widetilde{\mathcal M} \arrow[r] \arrow[d] &
\underbrace{\Gamma}_{N_{0,3}=\Gamma_{\mathbb Q}\times\mathfrak S_3}
\arrow[r,no head] \arrow[d,"\unsure{\chi}\times\operatorname{id}_{\mathfrak S_3}"] & 1 \\
1 \arrow[r,no head] &
\widehat M \arrow[r,no head] &
G_{\rho,\widetilde\sigma} \arrow[r] &
\widehat{\mathbb Z}^{*}\times\mathfrak S_3 \arrow[r,no head] & 1 .
\end{tikzcd}
\tag{282}
\]
}

Moreover against \unsure{all|all} expectation we found that the extension $G_{\rho,\widetilde\sigma}$ of $\widehat{\mathbb Z}^{*}\times\mathfrak S_3$ by $\widehat M$ was split, whence it follows that \deleted{that of} the extension $\widetilde{\mathcal M}$. of $\Gamma$ by $\widehat M$ is so too (\deleted{what \unsure{precedes|preceded}} \inserted{this last fact} is also a consequence of the existence of the Fermat tower, by applying to it the general reflections of the preceding pages). Better still, we have seen that there is a unique $\widehat M$-conjugacy class of splittings of the extension $G_{\rho,\widetilde\sigma}$, which induces a \uline{canonical} (but not unique!) $\widehat M$-conjugacy class of \deleted{$\unsure{\widehat M}$} splittings of $\widetilde{\mathcal M}$ over $\Gamma$. This therefore defines a canonical \deleted{isomorphism}\inserted{$\widehat M$-conjugacy} class of liftings of the action of $\Gamma$ \inserted{on $\mathcal X=U_{0,3\overline{\mathbb Q}}$} to $\mathcal X'_0$ -- or what amounts to the same thing, a torsor on $(U_{0,3\mathbb Q},\mathfrak S_3)$, with group a profinite group

\margin{\uline{N.B.} $H^1(\Gamma,\widehat M)$ is a very big group! -- But see \unsure{below}, cf. Theorem page~716}

\margin{126}

\begin{center}
\small
126\\
\texttt{https://grothendieck.umontpellier.fr}\\
Université de Montpellier, UMR CNRS 5149, Case courrier 051, Place Eugène Bataillon\\
34095 Montpellier cedex 5 - France
\end{center}